\section{Introduction}

Predictive models learned from data lie at the core of most practical AI systems. These models are typically trained to minimize the average prediction error. However, in high-stakes applications such as medicine, it is crucial not only to achieve low prediction error but also to quantify the uncertainty of individual predictions and effectively communicate this uncertainty to the end user.

Prediction uncertainty arises from two primary sources: aleatoric uncertainty and epistemic uncertainty~\cite{Hullermeier-Intro-ML2021}. Aleatoric uncertainty stems from the inherent randomness in the relationship between the inputs and the target outputs. It reflects noise in the data and is considered irreducible. In contrast, epistemic uncertainty arises due to the model being trained on a finite dataset instead of having access to the true data-generating distribution. Unlike aleatoric uncertainty, epistemic uncertainty is reducible with more data.

Prediction uncertainty can be leveraged to support decision-making in various ways. A widely used approach is reject-option prediction, where the model abstains from predicting when uncertainty is high, thus indicating an elevated risk of an incorrect outcome. Traditional reject-option predictors, as first proposed by~\cite{Chow-RejectOpt-TIT1970}, consider only aleatoric uncertainty, a setting we term aleatoric reject-option prediction. Existing methods for training these predictors often rely on empirical risk minimization or maximum likelihood estimation, assuming a training dataset large enough to disregard epistemic uncertainty~\cite{Hendrickx-MLwithRejectOpt-ML2024}.

In practical applications, training data is inherently finite, and consequently, ignoring epistemic uncertainty is often unjustified. Bayesian learning provides a principled methodology for naturally incorporating both aleatoric and epistemic uncertainties~\cite{Bishop-PRML-2006,Gelman-Bayesian-2013}. Its application to reject-option prediction enables the development of models we call Bayesian reject-option predictors. Such models are designed to abstain from making a prediction when the overall predictive uncertainty, encompassing both its aleatoric and epistemic sources, is deemed too high.

In this paper, we present a novel framework for constructing epistemic reject-option predictors that abstain from making predictions when epistemic uncertainty is high, i.e., when the training data is insufficient to support reliable decisions for a given input. To formalize this behavior, we extend the standard Bayesian learning framework by replacing the conventional objective of minimizing expected loss with the minimization of expected regret. Regret is defined as the performance gap between the learned predictor and the Bayes-optimal predictor, which has full knowledge of the true data-generating distribution. The model abstains when the regret for a given input exceeds a specified rejection cost.

The main contributions of this work are:
\begin{enumerate}
\item We propose a principled framework for constructing reject-option predictors that base their accept-or-reject decisions solely on epistemic uncertainty. To our knowledge, this is the first framework that enables predictors to systematically identify inputs for which available training data is insufficient to support reliable predictions.
%
\item Our framework provides a theoretical justification for entropy-based and variance-based measures of epistemic uncertainty commonly used in Bayesian neural networks~\cite{Depeweg-Decompos-ICML2018}. Specifically, we show that these widely used uncertainty measures correspond to the conditional regret, the quantity our optimal epistemic reject-option predictor uses to decide whether to accept or reject a prediction.
\end{enumerate}
